\par\Needspace{9\baselineskip}\bigskip
\pdfbookmark[3]{Exposition, 19}{pallas-II-article19}
\begin{center}19.\end{center}
\begingroup
\setlength{\lineskip}{4pt}\setlength{\lineskiplimit}{4pt}
\noindent\hspace*{\parindent}En appliquant ces méthodes d’intégration à nos variations des élémens, on aura, sauf les constantes que l’intégration introduit, les valeurs de \(i,\ \AscendingNode,\ \varphi,\ n,\ \tilde\omega,\ \epsilon\), pour les époques intermédiaires à celles pour lesquelles on a calculé les variations instantanées, et la valeur de \(\int n\mathrm{d}t\) pour ces époques mêmes, si on se sert des formules I, II; ou vice versâ, si on préfère les formules III, IV. On peut aussi pour l’intégration simple employer la formule III, et pour l’intégration double la formule II; on aura par là tous les élémens pour les mêmes époques. Au reste cela est assés indifférent, puisque ces valeurs des élémens ne doivent servir que pour en déduire, au moyen des méthodes connues d’interpolation, celles qui se rapportent aux époques des différentes observations. Le calcul de ces observations fournira, d’après des méthodes connues, les valeurs des constantes, qui sont au nombre de sept, mais qui se réduisent à six, puisque la constante \(\alpha\) qui résulte de \(\int\dfrac{\mathrm{d}\epsilon}{\mathrm{d}t}\mathrm{d}t\), et la constante \(\beta\) qui provient de \(\int n\mathrm{d}t\) ou \(\int\!\int\dfrac{\mathrm{d}n}{\mathrm{d}t}\mathrm{d}t^2\), se confondent en une seule.
\par\endgroup
